\title{AI existential risk probabilities are too unreliable to inform policy}
\begin{document}
\maketitle

\begin{abstract}
Von: Arvind and Sayash from AI as Normal Technology aisnakeoil@substack.com Betreff: AI existential risk probabilities are (still) too unreliable to inform policy Datum: 28. September 2026 um 13:20 An: johannes.meier@xigmbh.de Forwarded this email? Subscribe here for more AI existential risk probabilities are (still) too unreliable to inform policy How speculation gets laundered through pseudo- quantification ARVIND NARAYANAN AND SAYASH KAPOOR SEP 28 READ IN APP Over two years ago we wrote a detailed deconstruction of p(doom) and argued that its primary effect is to launder vague intuitions and fears through a facade of quantification. The essay has fresh relevance today as p(doom) rhetoric is driving public discourse and policy attention to an unprecedented degree. Unsurprisingly, today’s probability estimates of AI existential risk are no more rigorous than the ones from 2024. Hence this repost. We should be clear about what we’re not saying. We aren’t saying forecasters (or others who attempt to quantify AI x-risk) intend to mislead. Unfortunately, we think the effect is misleading nonetheless. In general, being quantitative about forecasts is good. In fact, Arvind has been serving as an advisor to the Forecasting Research Institute’s Longitudinal Expert AI Panel, an effort he thinks is extremely valuable. As the essay makes clear, the infeasibility of building useful quantitative models is specific to forecasts of AI x-risk, not forecasting in general. Our goal here is not 
\end{abstract}

\section{Recovered Text}
Von: Arvind and Sayash from AI as Normal Technology aisnakeoil@substack.com
Betreff: AI existential risk probabilities are (still) too unreliable to inform policy
Datum: 28. September 2026 um 13:20
An: johannes.meier@xigmbh.de
Forwarded this email? Subscribe here for more
AI existential risk probabilities
are (still) too unreliable to
inform policy
How speculation gets laundered through pseudo-
quantification
ARVIND NARAYANAN AND SAYASH KAPOOR
SEP 28
READ IN APP
Over two years ago we wrote a detailed deconstruction of p(doom)
and argued that its primary effect is to launder vague intuitions and
fears through a facade of quantification. The essay has fresh
relevance today as p(doom) rhetoric is driving public discourse and
policy attention to an unprecedented degree. Unsurprisingly, today’s
probability estimates of AI existential risk are no more rigorous than
the ones from 2024. Hence this repost.
We should be clear about what we’re not saying.
We aren’t saying forecasters (or others who attempt to quantify
AI x-risk) intend to mislead. Unfortunately, we think the effect is
misleading nonetheless.
In general, being quantitative about forecasts is good. In fact,
Arvind has been serving as an advisor to the Forecasting
Research Institute’s Longitudinal Expert AI Panel, an effort he
thinks is extremely valuable. As the essay makes clear, the
infeasibility of building useful quantitative models is specific to
forecasts of AI x-risk, not forecasting in general.
Our goal here is not to try to push AI safety advocates to stop
communicating in terms of probabilities (not that we’d succeed
even if we tried). Rather, our hope in publishing this is for those
at the receiving end of these probabilities — especially
policymakers — to recognize that unlike almost every other
quantitative forecast that we encounter, they are not the result
of a validated model or method.
That said, we think the whole p(doom) culture is emblematic of a
particular view of AI safety that centers existential risk and
superintelligence in a way that is actively counterproductive to a
broader conception of safety. In a few days we’ll publish a follow-up
essay making that argument and advocating for a “big tent” safety
movement.
The original essay appears below.
How seriously should governments take the threat of existential risk
from AI, given the lack of consensus among researchers? On the
one hand, existential risks (x-risks) are necessarily somewhat
AI existential risk probabilities
are too unreliable to inform policy
speculative: by the time there is concrete evidence, it may be too late.
On the other hand, governments must prioritize — after all, they don’t
worry too much about x-risk from alien invasions.
The AI safety community relies heavily on forecasting the probability
of human extinction due to AI (in a given timeframe) in order to
inform decision making and policy. An estimate of 10\% over a few
decades, for example, would obviously be high enough for the issue
to be a top priority for society.
Our central claim is that AI x-risk forecasts are far too unreliable to be
useful for policy, and in fact highly misleading.
If the two of us predicted an 80\% probability of aliens landing on
earth in the next ten years, would you take this possibility seriously?
Of course not. You would ask to see our evidence. As obvious as this
may seem, it seems to have been forgotten in the AI x-risk debate
that probabilities carry no authority by themselves. Probabilities are
usually derived from some grounded method, so we have a strong
cognitive bias to view quantified risk estimates as more valid than
qualitative ones. But it is possible for probabilities to be nothing more
than guesses. Keep this in mind throughout this essay (and more
broadly in the AI x-risk debate).
If we predicted odds for the Kentucky Derby, we don’t have to give
you a reason — you can take it or leave it. But if a policymaker takes
actions based on probabilities put forth by a forecaster, they had
better be able to explain those probabilities to the public (and that
explanation must in turn come from the forecaster). Justification is
essential to legitimacy of government and the exercise of power. A
Look behind the curtain
g
y
g
core principle of liberal democracy is that the state should not limit
people’s freedom based on controversial beliefs that reasonable
people can reject.
Explanation is especially important when the policies being
considered are costly, and even more so when those costs are
unevenly distributed among stakeholders. A good example is
restricting open releases of AI models. Can governments convince
people and companies who stand to benefit from open models that
they should make this sacrifice because of a speculative future risk?
The main aim of this essay is analyzing whether there is any
justification for any of the specific x-risk probability estimates that
have been cited in the policy debate. We have no objection to AI x-
risk forecasting as an academic activity, and forecasts may be helpful
to companies and other private decision makers. We only question
its use in the context of public policy.
There are basically only three known ways by which a forecaster can
try to convince a skeptic: inductive, deductive, and subjective
probability estimation. We consider each of these in the following
sections. All three require both parties to agree on some basic
assumptions about the world (which cannot themselves be proven).
The three approaches differ in terms of the empirical and logical
ways in which the probability estimate follows from that set of
assumptions.
Most risk estimates are inductive: they are based on past
observations. For example, insurers base their predictions of an
Inductive probability estimation is unreliable
due to the lack of a reference class
p ,
p
individual’s car accident risk on data from past accidents about
similar drivers. The set of observations used for probability
estimation is called a reference class. A suitable reference class for
car insurance might be the set of drivers who live in the same city. If
the analyst has more information about the individual, such as their
age or the type of car they drive, the reference class can be further
refined.
For existential risk from AI, there is no reference class, as it is an
event like no other. To be clear, this is a matter of degree, not kind.
There is never a clear “correct” reference class to use, and the
choice of a reference class in practice comes down to the analyst’s
intuition.
The accuracy of the forecasts depends on the degree of similarity
between the process that generates the event being forecast and the
process that generated the events in the reference class, which can
be seen as a spectrum. For predicting the outcome of a physical
system such as a coin toss, past experience is a highly reliable guide.
Next, for car accidents, risk estimates might vary by, say, 20\% based
on the past dataset used — good enough for insurance companies.
Further along the spectrum are geopolitical events, where the choice
of reference class gets even fuzzier. Forecasting expert Philip Tetlock
explains: “Grexit may have looked sui generis, because no country
had exited the Eurozone as of 2015, but it could also be viewed as
just another instance of a broad comparison class, such as
negotiation failures, or of a narrower class, such as a nation-states
withdrawing from international agreements or, narrower still, of forced
currency conversions.” He goes on to defend the idea that even
seeming Black Swan events like the collapse of the USSR or the
Arab Spring can be modeled as members of reference classes, and
that inductive reasoning is useful even for this kind of event.
In Tetlock’s spectrum, these events represent the “peak” of
uniqueness. When it comes to geopolitical events, that might be true.
But even those events are far less unique than extinction from AI.
Just look at the attempts to find reference classes for AI x-risk:
animal extinction (as a reference class for human extinction), past
global transformations such as the industrial revolution (as a
reference class for socioeconomic transformation from AI), or
accidents causing mass deaths (as a reference class for accidents
causing global catastrophe). Let’s get real. None of those tell us
anything about the possibility of developing superintelligent AI or
losing control over such AI, which are the central sources of
uncertainty for AI x-risk forecasting.
To summarize, human extinction due to AI is an outcome so far
removed from anything that has happened in the past that we cannot
use inductive methods to “predict” the odds. Of course, we can get
qualitative insights from past technical breakthroughs as well as past
catastrophic events, but AI risk is sufficiently different that
quantitative estimates lack the kind of justification needed for
legitimacy in policymaking.
D d
i
b bili
i
i
i
li bl
In Conan Doyle’s The Adventure of the Six Napoleons — spoiler
alert! — Sherlock Holmes announces before embarking on a
stakeout that the probability of catching the suspect is exactly two-
thirds. This seems bewildering — how can anything related to human
behavior be ascribed a mathematically precise probability?
It turns out that Holmes has deduced the underlying series of events
that gave rise to the suspect’s seemingly erratic observed behavior:
the suspect is methodically searching for a jewel that is known to be
hidden inside one of six busts of Napoleon owned by different people
in and around London. The details aren’t too important, but the key is
that neither the suspect nor the detectives know which of the six
busts it is in, and everything else about the suspect’s behavior is
(assumed to be) entirely predictable. Hence the precisely quantifiable
uncertainty.
The point is that if we have a model of the world that we can rely
upon, we can estimate risk through logical deduction, even without
relying on past observations. Of course, outside of fictional scenarios,
the world isn’t so neat, especially when we want to project far into the
future.
When it comes to x-risk, there is an interesting exception to the
general rule that we don’t have deductive models — asteroid impact.
A combination of inductive and deductive risk estimation does allow
us to estimate the probability of x-risk, only because we’re talking
about a purely physical system. Let’s take a minute to review how
this works, because it’s important to recognize that the methods are
not generalizable to other types of x-risk.
Deductive probability estimation is unreliable
due to the lack of theory
The key is being able to model the relationship between the size of
the asteroid (more precisely, the energy of impact) and the frequency
of impact. Since we have observed thousands of small impacts, we
can extrapolate to infer the frequency of large impacts that have
never been directly observed. We can also estimate the threshold
that would cause global catastrophe.¹
Figure: data on small asteroid impacts (illustrated on the
left) can be extrapolated to extinction-level impacts
(right).
With AI, the unknowns relate to technological progress and
governance rather than a physical system, so it isn’t clear how to
model it mathematically. Still, people have tried. For example, in
order to predict the computational requirements of a hypothetical AGI,
several works assume that an AI system would require roughly as
many computations as the human brain, and further make
assumptions about the number of computations required by the
human brain. These assumptions are far more tenuous than those
involved in asteroid modeling, and none of this even addresses the
loss-of-control question.
Subjective probabilities are feelings dressed
Without the reference classes or grounded theories, forecasts are
necessarily “subjective probabilities”, that is, guesses based on the
forecaster’s judgment. Unsurprisingly, these vary by orders of
magnitude.
Subjective probability estimation does not get around the need for
having either an inductive or a deductive basis for probability
estimates. It merely avoids the need for the forecaster to explain their
estimate. Explanation can be hard due to humans’ limited ability to
explain our intuitive reasoning, whether inductive, deductive, or a
combination thereof. Essentially, it allows the forecaster to say: “even
though I haven’t shown my methods, you can trust this estimate
because of my track record” (we explain in the next section why even
this breaks down for AI x-risk forecasting). But ultimately, lacking
either an inductive or a deductive basis, all that forecasters can do is
to make up a number, and those made-up numbers are all over the
place.
Subjective probabilities are feelings dressed
up as numbers
Consider the Existential Risk Persuasion Tournament (XPT)
conducted by the Forecasting Research Institute in late 2022, which
we think is the most elaborate and well-executed x-risk forecasting
exercise conducted to date. It involved various groups of forecasters,
including AI experts and forecasting experts (“superforecasters” in
the figure). For AI experts, the high end (75th percentile) of estimates
for AI extinction risk by 2100 is 12\%, the median estimate is 3\%, and
the low end (25th percentile) is 0.25\%. For forecasting experts, even
the high end (75th percentile) is only 1\%, the median is a mere
0.38\%, and the low end (25th percentile) is visually indistinguishable
from zero on the graph. In other words, the 75th percentile AI expert
forecast and the 25th percentile superforecaster forecast differ by at
least a factor of 100.
All of these estimates are from people who have deep expertise on
the topic and participated in a months-long tournament where they
tried to persuade each other! If this range of forecasts here isn’t
extreme enough, keep in mind that this whole exercise was
conducted by one group at one point in time. We might get different
numbers if the tournament were repeated today, if the questions
were framed differently, etc.
What’s most telling is to look at the rationales that forecasters
provided, which are extensively detailed in the report. They aren’t
using quantitative models, especially when thinking about the
likelihood of bad outcomes conditional on developing powerful AI.
For the most part, forecasters are engaging in the same kind of
speculation that everyday people do when they discuss
superintelligent AI. Maybe AI will take over critical systems through
superhuman persuasion of system operators. Maybe AI will seek to
lower global temperatures because it helps computers run faster, and
accidentally wipe out humanity. Or maybe AI will seek resources in
space rather than Earth, so we don’t need to be as worried. There’s
nothing wrong with such speculation. But we should be clear that
when it comes to AI x-risk, forecasters aren’t drawing on any special
knowledge, evidence, or models that make their hunches more
credible than yours or ours or anyone else’s.
The term superforecasting comes from Philip Tetlock’s 20 year study
of forecasting (he was also one of the organizers of the XPT).
Superforecasters tend to be trained in methods to improve forecasts
such as by integrating diverse information and by minimizing
psychological biases. These methods have been shown to be
effective in domains such as geopolitics. But no amount of training
will lead to good forecasts if there isn’t much useful evidence to draw
from.
Even if forecasters had credible quantitative models (they don’t), they
must account for “unknown unknowns”, that is, the possibility that the
model itself might be wrong. As noted x-risk philosopher Nick
Bostrom explains: “The uncertainty and error-proneness of our first-
order assessments of risk is itself something we must factor into our
all-things-considered probability assignments. This factor often
dominates in low-probability, high-consequence risks — especially
those involving poorly understood natural phenomena, complex
social dynamics, or new technology, or that are difficult to assess for
other reasons.”
This is a reasonable perspective, and AI x-risk forecasters do worry a
lot about uncertainty in risk assessment. But one consequence of
this is that for those who follow this principle, forecasts are
guaranteed to be guesses rather than the output of a model — after
all, no model can be used to estimate the probability that the model
itself is wrong, or what the risk would be if the model were wrong.
To recap, subjective AI-risk forecasts vary by orders of magnitude.
But if we can measure forecasters’ track records, maybe we can use
that to figure out which forecasters to trust. In contrast to the
previous two approaches for justifying risk estimates (inductive and
deductive), the forecaster doesn’t have to explain their estimate, but
instead justifies it based on their demonstrated skill at predicting
other outcomes in the past.
This has proved to be invaluable in the domain of geopolitical events,
and the forecasting community spends a lot of effort on skill
measurement. Many ways to evaluate forecasting skill exist, such as
calibration, the Brier score, the logarithmic score, or the Peer score
used on the forecasting competition website Metaculus.
But regardless of which method is used, when it comes to existential
risk, there are many barriers to assessing forecast skill for subjective
probabilities: the lack of a reference class, the low base rate, and the
long time horizon. Let’s look at each of these in turn.
Just as the reference class problem plagues the forecaster, it also
affects the evaluator. Let’s return to the alien landing example.
Consider a forecaster who has proved highly accurate at calling
Forecast skill cannot be measured when it
comes to unique or rare events
elections. Suppose this forecaster announces, without any evidence,
that aliens will land on Earth within a year. Despite the forecaster’s
demonstrated skill, this would not cause us to update our beliefs
about an alien landing, because it is too dissimilar to election
forecasting and we do not expect the forecaster’s skill to generalize.
Similarly, AI x-risk is so dissimilar to any past events that have been
forecast that there is no evidence of any forecaster’s skill at
estimating AI x-risk.
Even if we somehow do away with the reference class problem,
other problems remain — notably, the fact that extinction risks are
“tail risks”, or risks that result from rare events. Suppose forecaster A
says the probability of AI x-risk is 1\%, and forecaster B says it is 1 in
a million. Which forecast should we have more confidence in? We
could look at their track records. Say we find that forecaster A (who
has assigned a 1\% probability to AI x-risk) has a better track record.
It still doesn’t mean we should have more confidence in A’s forecast,
because skill evaluations are insensitive to overestimation of tail risks
In other words, it could be that A scores higher overall because A is
slightly better calibrated than B when it comes to everyday events
that have a substantial probability of occurring, but tends to
massively overestimate tail risks that occur rarely (for example, those
with a probability of 1 in a million) by orders of magnitude. No scoring
rule adequately penalizes this type of miscalibration.
Here’s a thought experiment to show why this is true. Suppose two
forecasters F and G forecast two different sets of events, and the
“true” probabilities of events in both sets are uniformly distributed
between 0 and 1. We assume, highly optimistically, that both F and G
know the true probability P[e] for every event e that they forecast. F
always outputs P[e], but G is slightly conservative, never predicting a
value less than 1\%. That is, G outputs P[e] if P[e] >= 1\%, otherwise
outputs 1\%.
By construction, F is the better forecaster. But would this be evident
from their track records? In other words, how many forecasts from
each would we have to evaluate so that there’s a 95\% chance that F
outscores G? With the logarithmic scoring rule, it turns out to be on
the order of a hundred million. With the Brier score, it is on the order
of a trillion.² We can quibble with the assumptions here but the point
is that if a forecaster systematically overestimates tail risks, it is
simply empirically undetectable.
The final barrier to assessing forecaster skill at predicting x-risk is
that long-term forecasts take too long to evaluate (and extinction
forecasts are of course impossible to evaluate). This can potentially
be overcome. Researchers have developed a method called
reciprocal scoring — where forecasters are rewarded based on how
well they predict each others’ forecasts — and validate it in some
real-world settings, such as predicting the effect of Covid-19 policies.
In these settings, reciprocal scoring yielded forecasts that are as
good as traditional scoring methods. Fair enough. But reciprocal
scoring is not a way around the reference class problem or the tail
risk problem.
Summary of our argument so far, showing why none of
the three forecasting methods can yield credible
estimates of AI x-risk
estimates of AI x risk.
To recap, inductive and deductive methods don’t work, subjective
forecasts are all over the place, and there’s no way to tell which
forecasts are more trustworthy.
So in an attempt to derive more reliable estimates that could
potentially inform policy, some researchers have turned to forecast
aggregation methods that combine the predictions of multiple
forecasters. A notable effort is the AI Impacts Survey on Progress in
AI, but it has been criticized for serious methodological limitations
including non-response bias. More importantly, it is unclear why
aggregation should improve forecast accuracy: after all, most
forecasters might share the same biases (and again, none of them
have any basis for a reliable forecast).
There are many reasons why forecasters might systematically
overestimate AI x-risk.³ The first is selection bias. Take AI
researchers: the belief that AI can change the world is one of the
main motivations for becoming an AI researcher. And once someone
enters this community, they are in an environment where that
message is constantly reinforced. And if one believes that this
technology is terrifyingly powerful, it is perfectly rational to think there
is a serious chance that its world-altering effects will be negative
rather than positive.
And in the AI safety subcommunity, which is a bit insular, the echo
chamber can be deafening. Claiming to have a high p(doom) (one’s
estimate of the probability of AI doom) seems to have become a way
There are many reasons why risk estimates
may be systematically inflated
to signal one’s identity and commitment to the cause.
There is a slightly different selection bias at play when it comes to
forecasting experts. The forecasting community has a strong overlap
with effective altruism and concerns about existential risk, especially
AI risk. This doesn’t mean that individual forecasters are biased. But
having a high p(doom) might make someone more inclined to take
up forecasting as an activity. So the community as a whole is likely
biased toward people with x-risk worries.
Forecasters are good at updating their beliefs in response to
evidence, but the problem is that unlike, say, asteroid impact risk,
there is little evidence that can change one’s beliefs one way or
another when it comes to AI x-risk, so we suspect that forecasts are
strongly influenced by the priors with which people enter the
community. The XPT report notes that “Few minds were changed
during the XPT, even among the most active participants, and
despite monetary incentives for persuading others.” In a follow-up
study, they found that many of the disagreements were due to
fundamental worldview differences that go beyond AI.
To reemphasize, our points about bias are specific to AI x-risk. If
there were a community of election forecasters who were
systematically biased (say, toward incumbents), this would become
obvious after a few elections when comparing predictions with reality.
But with AI x-risk, as we showed in the previous section, skill
evaluation is insensitive to overestimation of tail risks.
Interestingly, skill evaluation is extremely sensitive to underestimation
of tail risks: if you assign a probability of 0 for a rare event that
actually ends up occurring, you incur an infinite penalty under the
logarithmic scoring rule from which you can never recover
logarithmic scoring rule, from which you can never recover
regardless of how well you predicted other events. This is considered
one of the main benefits of the logarithmic score and is the reason it
is adopted by Metaculus.
Now consider a forecaster who doesn’t have a precise estimate —
and surely no forecaster has a precise estimate for something with
so many axes of uncertainty as AI x-risk. Given the asymmetric
penalties, the rational thing to do is to go with the higher end of their
range of estimates.⁴
In any case, it’s not clear what forecasters actually report when their
estimates are highly uncertain. Maybe they don’t respond to the
incentives of the scoring function. After all, long-term forecasts won’t
be resolved anytime soon. And recall that in the case of the XPT, the
incentive is actually to predict each others’ forecasts to get around
the problem of long time horizons. The reciprocal scoring paper
argues that this will incentivize forecasters to submit their true, high-
effort estimates, and considers various objections to this claim. Their
defense of the method rests on two key assumptions: that by
exerting more effort forecasters can get closer to the true estimate,
and that they have no better way to predict what other forecasters
will do.
What if these assumptions are not satisfied? As we have argued
throughout this post, with AI x-risk, we shouldn’t expect evidence to
change forecasters’ prior beliefs, so the first assumption is dubious.
And now that one iteration of the XPT has concluded, the published
median estimates from that tournament serve as a powerful anchor
(a “focal point” in game theory). It is possible that in the future,
forecasters with reciprocal scoring incentives will use existing median
forecasts as a starting point, only making minor adjustments to
g p
,
y
g
j
account for new information that has become available since the last
tournament. The range of estimates might narrow as existing
estimates serve as anchors for future estimates. All that is a
roundabout way to say: the less actual evidence there is to draw
upon, the more the risk of groupthink.
For what it’s worth, here are the median estimates from the XPT of
both extinction risk and sub-extinction catastrophic risks from AI:
To reiterate, our view is that we shouldn’t take any of these numbers
too seriously. They are a reflection of how much different samples of
participants fret about AI than anything else.
As before, the estimates from forecasting experts (superforecasters)
and AI experts differ by an order of magnitude or more. To the extent
that we put any stock into these estimates, it should be the
forecasting experts’ rather than the AI experts’ estimates. One
important insight from past research is that domain experts perform
worse than forecasting experts who have training in integrating
diverse information and by minimizing psychological biases. Still, as
we said above even their forecasts may be vast overestimates and
Beware Pascal’s wager: the dangers of utility
maximization
we said above, even their forecasts may be vast overestimates, and
we just can’t know for sure.
So what’s the big deal? So what if policymakers believe the risk over
a certain timeframe is 1\% instead of 0.01\%? It seems pretty low in
either case!
It depends on what they do with those probabilities. Most often, these
estimates are merely a way to signal the fact that some group of
experts thinks the risk is significant. If that’s all they are, so be it. But
it’s not clear that all this elaborate effort at quantification is even
helpful for this signaling purpose, given that different people interpret
the same numbers wildly differently.
For example, Federal Trade Commission chair Lina Khan said her
views on the matter were techno-optimistic since her p(doom) was
only 15\%, which left experts bewildered. (For what it’s worth, that
number is about a thousandfold higher than what we would be
comfortable labeling techno-optimist.) It takes a lot of quantitative
training to be able to mentally process very small or very large
numbers correctly in decision making, and not simply bucket them
into categories like “insignificantly small”. Most people are not trained
this way.
In short, what seems to be happening is that experts’ vague intuitions
and fears are being translated into pseudo-precise numbers, and
then translated back into vague intuitions and fears by policymakers.
Let’s just cut the charade of quantification! The Center for AI Safety’s
Statement on AI Risk was admirably blunt in this regard (of course,
we strongly disagree with its substance).
A principled, quantitative way to use probabilities in decision making
is utility maximization through cost-benefit analysis. The idea is
simple: if we consider an outcome to have a subjective value, or
utility, of U (which can be positive or negative), and it has, say, a 10\%
probability of occurring, we can act as if it is certain to occur and has
a value of 0.1 * U. We can then add up the costs and benefits for
each option available to us, and choose the one that maximizes
costs minus benefits (the “expected utility”).
This is where things get really problematic. First, some people might
consider extinction to have an unfathomably large negative value,
because it precludes the existence of all the human lives, physical or
simulated, that might ever be born in the future. The logical
conclusion is that x-risk should be everyone’s top priority all the time!
It is reminiscent of Pascal’s wager, the argument that it is rational
believe in God because even if there is an infinitesimally small
chance that God exists, the cost of non-belief is infinite (an eternity in
hell as opposed to eternal happiness), and hence so is the expected
utility. Fortunately, policymakers don’t give too much credence to
decision making frameworks involving infinities. But the idea has
taken a powerful hold of the AI safety community and drives some
people’s conviction that AI x-risk should be society’s top priority.
Even if we limit ourselves to catastrophic but not existential risks, we
are talking about billions of lives on the line, so the expected cost of
even a 1\% risk is so high that the policy implications are drastic —
governments should increase spending on AI x-risk mitigation by
orders of magnitude and consider draconian measures such as
stopping AI development. This is why it is so vital to understand that
these estimates are not backed by any methodology. It would be
incredibly unwise to make world-changing policy decisions based on
so little evidence.
Is there a role for forecasting in AI policy? We think yes — just not
forecasting existential risk. Forecasting AI milestones, such as
performance on certain capability benchmarks or economic impacts,
is more achievable and meaningful. If a forecaster has demonstrated
skill in predicting when various AI milestones would be reached, it
does give us evidence that they will do well in the future. We are no
longer talking about unique or rare events. And when considering
lower-stakes policy interventions — preparing for potential economic
disruption rather than staving off killer robots — it is less critical that
forecasts be justified to the satisfaction of every reasonable person.
The forecasting community devotes a lot of energy to milestone
forecasting. On Metaculus, the question “Will there be Human-
machine intelligence parity before 2040?” has an aggregate
prediction of 96\% based on over 1,300 forecasters. That’s
remarkable! If we agreed with this forecast, we would be in favor of
the position that managing the safe transition to AGI should be a
global priority. Why don’t we?
The answer is in the fine print. There is no consensus on the
definition of a fuzzy concept such as AGI. Even if we fix a definition,
determining whether it has been achieved can be hard or impossible.
For effective forecasting, it is extremely important to avoid
ambiguous outcomes. The way the forecasting community gets
around this is by defining it in terms of relatively narrow skills, such
as exam performance.
The Metaculus intelligence parity question is defined in terms of the
Forecasts of milestones suffer from outcome
ambiguity
performance on graduate exams in math, physics, and computer
science. Based on this definition, we do agree with the forecast of
96\%. But we think the definition is so watered down that it doesn’t
mean much for policy. Forget existential risk — as we’ve written
before, AI performance on exams has so little construct validity that it
doesn’t even let us predict whether AI will replace workers.
Other benchmarks aren’t much better. In short, forecasting AI
capability timelines is tricky because of the huge gap between
benchmarks and real-world implications. Fortunately, better
benchmarks reflecting consequential real-world tasks are being
developed. In addition to benchmarks, we need naturalistic
evaluation, even if it is more costly. One type of naturalistic
evaluation is to measure how people perform their jobs differently
with AI assistance. Directly forecasting economic, social, or political
impacts — such as labor market transformation or AI-related
spending by militaries — could be even more useful, although harder
to unambiguously define and measure.
The responsibility for avoiding misuses of probability in policy lies
with policymakers. We are not calling for forecasters to stop
publishing forecasts in order to “protect” policymakers from being
misled. That said, we think forecasts should be accompanied by a
clear explanation of the process used and evidence considered. This
would allow policymakers to make informed decisions about whether
the justification presented meets the threshold that they are
comfortable with. The XPT is a good example of transparency, as is
this paper (though it is not about x-risk). On the other hand, simply
surveying a bunch of researchers and presenting aggregate numbers
is misinformative and should be ignored by policymakers
Concluding thoughts
is misinformative and should be ignored by policymakers.
So what should governments do about AI x-risk? Our view isn’t that
they should do nothing. But they should reject the kind of policies
that might seem compelling if we view x-risk as urgent and serious,
notably: restricting AI development. As we’ll argue in a future essay
in this series, not only are such policies unnecessary, they are likely
to increase x-risk. Instead, governments should adopt policies that
are compatible with a range of possible estimates of AI risk, and are
on balance helpful even if the risk is negligible. Fortunately, such
policies exist. Governments should also change policymaking
processes so that they are more responsive to new evidence. More
on all that soon.
The XPT report is titled Forecasting Existential Risks: Evidence
from a Long-Run Forecasting Tournament.
Tetlock and Gardner’s book Superforecasting summarizes
research by Tetlock, Barbara Mellers, and others.
Scott Alexander refutes the claim that there is something wrong
in principle with ascribing probabilities to unique events (we
largely agree). Our argument differs in two key ways from the
position he addresses. We aren’t talking about forecasting in
general; just its application to policymaking. And we don’t object
to it in principle. Our argument is empirical: AI x-risk forecasts
are extremely unreliable and lack justification. Theoretically this
could change in the future, though we aren’t holding our breath.
A paper by Friedman and Zuckhauser explains why probabilities
aren’t the whole story: two forecasts that have the same risk
estimate might have very different implications for policymakers.
In our view AI x risk forecasts fare poorly on two of the three
Further reading
In our view, AI x-risk forecasts fare poorly on two of the three
dimensions of confidence: a sound basis in evidence and a
narrow range of reasonable opinion.
A paper by our Princeton CITP colleagues led by Shazeda
Ahmed explains the epistemic culture of AI safety, which
consists of “cohesive, interwoven social structures of
knowledge-production and community-building”. It helps
understand why practices such as forecasting have become
pillars of how the AI safety community forms its beliefs, as
opposed to the broader scientific community that centers
practices such as peer review. Of course, we shouldn’t reject a
view just because it doesn’t conform to scientific orthodoxy. But
at the same time, we shouldn’t give any deference to the self-
styled AI safety community’s views on AI safety. It is important
to understand that the median member of the AI safety
community holds one particular stance on AI safety — a stance
that is highly contested and in our view rather alarmist.
We have written extensively about evidence-based AI safety.
Our best-known work includes the essay AI safety is not a
model property and the paper titled On the Societal Impact of
Open Foundation Models which was the result of a large
collaboration.
Acknowledgements. We are grateful to Benjamin Edelman, Ezra
Karger, Matt Salganik, and Ollie Stephenson for feedback on a draft.
This series of essays is based on an upcoming paper that benefited
from feedback from many people, including Seth Lazar and members
of the MINT lab at Australian National University, students in the
Limits to Prediction course at Princeton, Shazeda Ahmed, and
Zachary Siegel.
1
Specifically, by using data from nuclear weapons tests and
making a few physics-informed assumptions, we can calculate
what it would take to kick up enough dust to darken the skies for
a prolonged period and lead to a collapse of global agriculture.
2
Suppose G, the conservative forecaster, has a floor of ϵ for their
forecasts (in our example, 0.01). There is only an ϵ fraction of
events where the difference between the two forecasters is even
relevant. Even limiting ourselves to the set of events with a true
probability less than ϵ, some calculation shows that F’s expected
log score is better than G’s by a tiny amount — O(ϵ), and with
the Brier score, O(ϵ^2). So when looking at all events, the
differences in expected scores are O(ϵ^2) and O(ϵ^3)
respectively. Meanwhile, the variance in their scores comes out
to about 0.1 for both scoring functions. To be able to confidently
assert that the means of the two forecaster’s scores are different,
we need the number of events N to be large enough that the
standard deviation of the difference in mean scores, which is
sqrt(0.1 / N), to be much less than the expected difference. So
we need N ~ O(1/ϵ^4) for the log score and N ~ O(1/ϵ^6) for the
Brier score.
3
Of course, there are also reasons why forecasters might
underestimate the risk. We find those reasons less persuasive,
but more importantly, we consider them out of scope for this
discussion. The reason is simple: the advocacy for governments
to implement freedom-restricting policies is being justified by
high x-risk estimates. The burden of evidence is thus asymmetric
Those who call for such policies must justify their probability
estimates, including responding to concerns about upward
biases. If, in some strange future world, some people advocate
for restrictive policies based on low x-risk estimates, it will
become important to subject potential underestimation to the
same scrutiny.
4
In the theory of forecasting, the logarithmic score and other so-
called proper scoring rules are mathematically guaranteed to
elicit the forecaster’s “true belief”. But the definition of a proper
scoring rule assumes that the true belief is a single value,
whereas it tends to be a wide range. In such a scenario, the
forecaster is incentivized to report the mean of their distribution,
but what we intuitively mean by the “true belief” might better
correspond to the median.
LIKE
COMMENT
RESTACK
© 2026 Sayash Kapoor and Arvind Narayanan
Princeton University, NJ 08544
Unsubscribe
\end{document}
